\SourcePage{213}
\SGAfrontchapter{原书记号索引}
\noindent\textit{Index des notations}

\noindent 以下保留原索引的讲次、节号及所列原页码；每项附有通往本译本相应正文原页起点的链接。

\begin{longtable}{@{}p{.50\textwidth}p{.46\textwidth}@{}}
\textbf{记号} & \textbf{原书出处及正文链接}\\\hline
$\Gamma_Z,\ \underline{\Gamma}_Z$ & I 1；\SGAsourceRef{13}\\
$i_!,\ i^!$ & I 1；\SGAsourceRef{13}\\
$\mathbf Z_{Z,X}$ & I 1.6；\SGAsourceRef{16}\\
$\mathrm H^i_Z(X,F),\ \underline{\mathrm H}^i_Z(F)$ & I 2；\SGAsourceRef{18}\\
$\mathrm H^i_J(M),\ \mathrm H^i((f),M),\ \mathrm H^i(M)$ & IV 5.4；\SGAsourceRef{53}；V 2（原第 50 页）；\SGAsourceRef{58}\\
$\operatorname{Ass}M,\ \operatorname{Supp}M,\ \operatorname{Ann}M,\ \mathrm r(\mathfrak a)$ & III 1.1；\SGAsourceRef{29}\\
$\operatorname{prof}_I(M),\ \operatorname{prof}_Y(F)$ & III 2.3、2.8；\SGAsourceRef{32}、\SGAsourceRef{33}\\
$\underline{\mathrm{Ab}}$ & IV 1（原第 33 页）；\SGAsourceRef{41}\\
$\underline{\operatorname{Hom}}^\bullet(F^\bullet,G^\bullet)$ & V 1.1；\SGAsourceRef{55}\\
$\operatorname{Ext}^i_Z(X;F,G),\ \underline{\operatorname{Ext}}^i_Z(F,G)$ & VI 1.1；\SGAsourceRef{65}\\
$\mathbf{Et}(X),\ \mathbf L(Z)$ & X（原第 89 页）；\SGAsourceRef{97}\\
$\operatorname{Lef}(X,Y),\ \operatorname{Leff}(X,Y)$ & X 2（原第 90 页）；\SGAsourceRef{98}\\
$\underline{\mathrm P}(Z),\ \operatorname{Pic}(Z)$ & XI（原第 99 页）；\SGAsourceRef{107}\\
$\Pi_i^x(X)$ & XIII（原第 146 页）；\SGAsourceRef{154}\\
$\mathrm D^+(X)$ & XIV 1.0；\SGAsourceRef{167}、\SGAsourceRef{168}\\
$\operatorname{prof}_Y(F)\geq n$ & XIV 1.2；\SGAsourceRef{170}\\
$\operatorname{prof}^{\mathbf L}_Y(X)\geq n$ & XIV 1.2；\SGAsourceRef{171}\\
$\operatorname{prof\,\text{ét}}_Y(X)$ & XIV 1.2；\SGAsourceRef{171}\\
$\mathbf L$ & XIV 1.2；\SGAsourceRef{171}\\
$\operatorname{prof}_x(F)\geq n$ & XIV 1.7；\SGAsourceRef{173}\\
$\operatorname{prof}^{\mathbf L}_x(X)\geq n$ & XIV 1.7；\SGAsourceRef{173}\\
$\operatorname{prof\,hop}^{\mathbf L}_x(X)\geq 3$ & XIV 1.7；\SGAsourceRef{173}\\
$\operatorname{prof\,\text{ét}}_x(X)$ & XIV 1.7；\SGAsourceRef{173}\\
$\delta_t(x),\ \delta(n)$ & XIV 2.1；\SGAsourceRef{185}\\
$\mathrm H^i_{Z!}(X/S,F)$ & XIV 4.0；\SGAsourceRef{195}\\
\end{longtable}

\SourcePage{214}
% 原 PDF 第 214 页为空白页。

\SourcePage{215}
\SGAfrontchapter{原书术语索引}
\noindent\textit{Index terminologique}

\noindent 保留原书的法文词头与英文对照；出处中的页码仍为原索引所列页码。

\begin{description}[style=nextline,leftmargin=0pt,labelindent=0pt,font=\normalfont\bfseries]
\item[Auslander-Buchsbaum（théorème de -）]
Auslander--Buchsbaum 定理。XI 3.13；\SGAsourceRef{113}。

\item[comparaison（théorème de -）]
比较定理（comparison theorem）。IX 1.1、XII 2.1；\SGAsourceRef{88}、\SGAsourceRef{122}。

\item[dualité locale（théorème de -）]
局部对偶定理（local duality theorem）。V 2.1；\SGAsourceRef{58}。

\item[dualité projective（théorème de -）]
射影对偶定理（projective duality theorem）。XII 1.1；\SGAsourceRef{117}。

\item[enveloppe injective]
内射包络（injective envelope）。IV（原第 41 页）；\SGAsourceRef{49}。

\item[excision（théorème d’ -）]
切除定理（excision theorem）。I 2.2、VI 1.3；\SGAsourceRef{18}、\SGAsourceRef{65}。

\item[existence（théorème d’ -）]
存在定理（existence theorem）。IX 2.1、XII 3.1；\SGAsourceRef{93}、\SGAsourceRef{125}。

\item[extension essentielle]
本质扩张（essential extension）。IV（原第 41 页）；\SGAsourceRef{49}。

\item[finitude（théorème de -）]
有限性定理（finiteness theorem）。VIII 2.1、XII 1.5；\SGAsourceRef{79}、\SGAsourceRef{119}。

\item[foncteur dualisant]
对偶化函子（dualizing functor）。IV 4.1、4.2；\SGAsourceRef{47}。

\item[forme résidu]
留数形式（residue form）。IV 5.5；\SGAsourceRef{53}。

\item[géométriquement factoriel, resp. géom. parafactoriel（anneau local -）]
几何阶乘局部环，分别地，几何副阶乘局部环（geometrically factorial, resp. geom. parafactorial local ring）。XIII（原列第 20、24 页）\footnote{译者注：原索引确实印作“XIII（p. 20, 24）”，此处照录。本版第 XIII 讲讨论这两个概念的段落见原印刷第 150、152 页，即 \SGAsourceRef{158}、\SGAsourceRef{160}。}；\SGAsourceRef{158}、\SGAsourceRef{160}。

\item[groupes d’homotopie locale]
局部同伦群（local homotopy groups）。XIII（原第 146 页）；\SGAsourceRef{154}。

\item[Gysin（homomorphisme de -）]
Gysin 同态（Gysin homomorphism）。I（原第 14 页）；\SGAsourceRef{22}。

\item[Hartshorne（théorème de -）]
Hartshorne 定理。III 3.6；\SGAsourceRef{36}。

\item[Lefschetz（condition de -）]
Lefschetz 条件（Lefschetz condition）。X 2（原第 90 页）；\SGAsourceRef{98}。

\item[Lefschetz（condition de - effective）]
有效 Lefschetz 条件（effective Lefschetz condition）。X 2（原第 90 页）；\SGAsourceRef{98}。

\item[Lefschetz（théorème de - affine）]
仿射 Lefschetz 定理（affine Lefschetz theorem）。XIV 3.1；\SGAsourceRef{189}。

\item[module dualisant]
对偶化模（dualizing module）。IV 4.1；\SGAsourceRef{47}。

\item[orthogonal d’un sous-module]
子模的正交补（orthogonal of a submodule）。IV 5；\SGAsourceRef{51}。

\item[parafactoriel（couple - de préschémas）]
预概形的副阶乘对（parafactorial pair of preschemes）。XI 3.1；\SGAsourceRef{109}。

\item[parafactoriel（anneau local -）]
副阶乘局部环（parafactorial local ring）。XI 3.2；\SGAsourceRef{109}。

\item[profondeur]
深度（depth）。III 2.3；\SGAsourceRef{32}。

\item[profondeur étale]
平展深度（etale depth）。XIV 1.2；\SGAsourceRef{170}、\SGAsourceRef{171}。

\item[profondeur géométrique d’un anneau local noethérien]
诺特局部环的几何深度（geometrical depth of a noetherian local ring）。XIV 5.3；\SGAsourceRef{208}。

\item[profondeur homotopique]
同伦深度（homotopical depth）。XIII 6 定义 1（原列第 154 页）；\SGAsourceRef{163}。\footnote{译者注：原索引列第 154 页；定义 1 实际位于本版原印刷第 155 页，此处的正文链接指向该定义。}

\item[profondeur homotopique（pour $\mathbf L$）]
关于 $\mathbf L$ 的同伦深度（homotopical depth with respect to $\mathbf L$）。XIV 1.2；\SGAsourceRef{171}。

\item[profondeur homotopique rectifiée]
校正同伦深度（rectified homotopical depth）。XIII 定义 2（原第 155 页）；\SGAsourceRef{163}。

\SourcePage{216}
\item[pur（anneau local -）]
纯局部环（pure local ring）。X 3.2；\SGAsourceRef{102}。

\item[pure（couple - de préschémas）]
预概形的纯对（pure pair of preschemes）。X 3.1；\SGAsourceRef{102}。

\item[pureté（théorème de - de Zariski-Nagata）]
Zariski--Nagata 纯性定理（purity theorem of Zariski-Nagata）。X 3.4；\SGAsourceRef{103}。

\item[pureté cohomologique（théorème de -）]
上同调纯性定理（cohomological purity theorem）。XIV 1.11；\SGAsourceRef{177}。

\item[régulier（$M$-régulier）]
正则（$M$-正则）（$M$-regular）。2.1；\SGAsourceRef{30}。\footnote{译者注：原索引此项只列“2.1”，未列讲次；第 III 讲引理 2.1 的正文给出 $M$-正则的定义，上列链接指向该页。}

\item[Samuel（conjecture de -）]
Samuel 猜想。XI 3.14；\SGAsourceRef{113}。

\item[semi-pureté cohomologique（théorème de -）]
上同调半纯性定理（cohomological semi-purity theorem）。XIV 1.10；\SGAsourceRef{177}。

\item[strictement local（anneau -）]
严格局部环（strictly local ring）。XIII 6（原第 151 页）；\SGAsourceRef{159}。

\item[Zariski-Nagata（théorème de pureté de -）]
Zariski--Nagata 纯性定理（purity theorem of Zariski-Nagata）。X 3.4；\SGAsourceRef{103}。
\end{description}
